\chapter{Geometry from natural language (LLM interface)}\label{ch:llm}

\NF{} can build geometry from requests in natural language, such as ``a box of
$20\times10\times5$\,mm with a bore of diameter 4 through its centre''. A local large language
model (LLM) served by \href{https://ollama.com}{Ollama} answers the request by calling geometry
\emph{tools}. The subpackage \code{nurbforge.llm} provides:

\begin{itemize}
  \item a minimal HTTP client for Ollama (\apiref{llm.OllamaClient});
  \item a whitelisted set of 40 geometry tools with JSON schemas (\code{nurbforge.llm.tools});
  \item the tool-calling loop with automatic verification (\apiref{llm.GeometryAgent});
  \item a measurement and checking layer (\code{nurbforge.llm.verify});
  \item a constraint solver for dimensions given only indirectly
    (\code{nurbforge.llm.constraints});
  \item an interactive command line, \code{nurbforge-llm}.
\end{itemize}

Everything runs locally. No data leaves the machine, and no API key is needed.

\section{Safety model}\label{sec:llm-safety}

The model \textbf{never executes code}. It can only call the tools of a fixed list. Each tool
has a JSON schema, and each operates on a named object store, the \apiref{llm.Workspace}. This
has three consequences:

\begin{itemize}
  \item \textbf{Bounded effects.} A tool can create, copy or measure objects, and write or read
    STEP and mesh files at paths the model chooses. Nothing else is possible. In particular
    there is no \code{eval}, no shell and no Python import.
  \item \textbf{Errors are text.} An invalid call, such as an unknown object, a wrong argument
    or degenerate geometry, does not raise into the host program. It is returned to the model
    as a message beginning with \code{ERROR:}, so the model can correct itself.
  \item \textbf{Arithmetic without Python.} Even the constraint solver's expressions
    (\cref{sec:llm-solve}) are evaluated by a small whitelist interpreter, not by Python.
\end{itemize}

This makes the interface usable in automated pipelines. The worst a confused model can do is
create wrong geometry, and the verification layer is there to catch that.

\begin{warning}[File paths]
The tools \toolref{export_step} and \toolref{import_step} read and write the file paths that
the model gives. Run the agent in a working directory where this is acceptable, or remove the
tools from \code{TOOLS} if your application must not touch files.
\end{warning}

\section{Setting up Ollama}

\begin{lstlisting}[style=shell]
# install Ollama from https://ollama.com, then:
ollama serve &                  # starts the server on http://localhost:11434
ollama pull qwen2.5:14b         # any model with tool support
nurbforge-llm --model qwen2.5:14b
\end{lstlisting}

The model must support \emph{tool calling}. Suitable model families include \code{qwen2.5}
(\code{qwen2.5:14b} is the default), \code{llama3.1} and \code{mistral-nemo}. As a rule, larger
models choose operations and arguments more reliably. The error messages and verification rounds
help smaller models, but cannot fully make up for them. Requests are sent with temperature 0,
so answers are largely reproducible.

Two environment variables set the defaults:
\begin{itemize}
  \item \code{OLLAMA_HOST} (default \code{http://localhost:11434});
  \item \code{NURBFORGE_LLM_MODEL} (default \code{qwen2.5:14b}).
\end{itemize}

\section{Using the agent from Python}

\begin{lstlisting}[style=python]
from nurbforge.llm import GeometryAgent, OllamaClient

agent = GeometryAgent(OllamaClient(host="http://localhost:11434", model="qwen2.5:14b"))
print(agent.ask("Revolve a semicircle of radius 5 about z into a sphere named ball"))
ball = agent.workspace.get("ball")       # an ordinary nurbforge object
\end{lstlisting}

\subsection{\texttt{OllamaClient}}

\code{OllamaClient(host=..., model=..., timeout=300.0)} is a thin wrapper around Ollama's REST
API. It uses only the standard library (\code{urllib}).
\begin{itemize}
  \item A host without a scheme gets \code{http://} prepended.
  \item \code{list_models()} returns the installed models (\code{/api/tags}).
  \item \code{chat(messages, tools=None, temperature=0.0)} performs one non-streaming chat turn
    (\code{/api/chat}) and returns the assistant message as a dictionary.
  \item Network problems and server errors raise \apiref{llm.OllamaError}, a
    \code{RuntimeError}.
\end{itemize}
Any object with a compatible \code{chat} method can replace the client. This is how the tests
work, and how the scripted example below runs without a model.

\subsection{\texttt{GeometryAgent}}

\begin{longtable}{@{}>{\raggedright\arraybackslash}p{0.3\textwidth}p{0.64\textwidth}@{}}
\toprule
\textbf{Member} & \textbf{Meaning} \\
\midrule\endhead
\code{GeometryAgent(client, workspace, max_steps, verify_rounds)} &
  All arguments are optional (defaults \code{None}, \code{None}, 20, 1). \code{client} defaults to \code{OllamaClient()}, and \code{workspace} to a new, empty
  \code{Workspace}. Pass an existing workspace to let the model work on objects you created in
  Python. \code{max_steps} limits the number of model turns per request phase. \\
\code{ask(prompt)} & Sends a request, executes the tool calls and runs the verification rounds
  (\cref{sec:llm-verify}). Returns the model's final answer. If the step limit is reached, the
  answer is \code{"Stopped: maximum number of tool steps reached."}. \\
\code{workspace} & The \code{Workspace} holding all objects. \\
\code{messages} & The complete conversation: system prompt, requests, tool calls and results.
  It persists across \code{ask} calls, so follow-up requests can refer to earlier objects. \\
\code{log} & List of \code{(tool, arguments, result)} for every executed call. \\
\code{verification} & The last verification report sent to the model, or \code{None}. \\
\bottomrule
\end{longtable}

The system prompt tells the model:
\begin{itemize}
  \item to build geometry only with the tools;
  \item that angles are given in degrees;
  \item to name objects descriptively and reuse existing names;
  \item to verify results with \code{inspect} and \code{check};
  \item to use \code{solve} instead of guessing when a dimension is given only indirectly.
\end{itemize}

\section{The tools}

\subsection{Overview}

\begin{longtable}{@{}p{0.22\textwidth}p{0.72\textwidth}@{}}
\toprule
\textbf{Group} & \textbf{Tools} \\
\midrule\endhead
Curves & \toolref{make_line}, \toolref{make_circle}, \toolref{make_arc},
  \toolref{make_polygon}, \toolref{make_helix}, \toolref{round_corners},
  \toolref{interpolate_points} \\
Surfaces & \toolref{make_rectangle}, \toolref{revolve}, \toolref{extrude}, \toolref{ruled},
  \toolref{loft}, \toolref{sweep}, \toolref{coons}, \toolref{asphere} \\
Modification & \toolref{extend}, \toolref{offset}, \toolref{fillet}, \toolref{trim_circle},
  \toolref{split}, \toolref{section} \\
Copies & \toolref{translate}, \toolref{rotate}, \toolref{mirror}, \toolref{linear_array},
  \toolref{polar_array} \\
Solids & \toolref{sew}, \toolref{solid} (box, sphere, cylinder, cone, torus),
  \toolref{boolean}, \toolref{mass_properties} \\
Queries & \toolref{list_objects}, \toolref{evaluate}, \toolref{closest_point},
  \toolref{measure}, \toolref{intersect} \\
Verification & \toolref{inspect}, \toolref{check}, \toolref{solve} \\
Files & \toolref{export_step}, \toolref{import_step} \\
\bottomrule
\end{longtable}

\Cref{app:llmtools} lists every tool with its description and all parameters, generated from
the schemas the model actually receives.

\subsection{Conventions at the tool level}

\begin{itemize}
  \item \textbf{Angles are in degrees.} Models handle degrees much more reliably than radians.
    The Python and C++ API use radians, and the tools convert.
  \item \textbf{Points and vectors} are lists of three numbers.
  \item \textbf{Objects are referred to by name.} A tool that creates something takes a
    \code{name}, and the object is stored under it, replacing an existing object of that name.
    Tools that create several objects append suffixes: \code{name_0} \dots{} for arrays,
    \code{name_1} \dots{} for \code{split}, \code{section} and \code{import_step}.
  \item \textbf{Results are short texts}, for example
    \code{stored part: closed shell, volume 20.8584}. They are designed to give the model the
    information it needs for the next step.
\end{itemize}

\section{Using the tools without a model}

The tool layer is useful on its own, for scripting, for testing, and for reproducing what a
model did. \code{call_tool(workspace, name, arguments)} executes exactly what the agent would
execute for a model's tool call, and \code{TOOLS} contains the tool definitions with their
schemas.

\example{16_llm_tools.py}

\begin{lstlisting}[style=text]
40 tools, e.g. make_line, make_circle, make_arc, make_polygon, make_helix, round_corners ...
make_circle - Create a full circle (exact NURBS).
  parameters: {"type": "number"} required: ['name', 'center', 'normal', 'radius']
stored box: closed shell, volume 24
stored rod: closed shell, volume 4.71239
stored part: closed shell, volume 20.8584
stored a: <NurbsCurve degree=2 n_ctrl=3 rational>
length=3.141592654
ERROR: KeyError: "unknown object 'nothing'. Known: ['a', 'box', 'part', 'rod']"
ERROR: unknown tool 'no_such_tool'
part: type=shell, faces=7, edges=17, free_edges=0, non_manifold_edges=0, closed=True,
  consistently_oriented=True, area=62.99557432, volume=20.85840737,
  bbox={min: [0, 0, -7.090971193e-17], max: [2, 3, 4]}
PASS: volume = 24 (expected 24)
FAIL: bbox.max = [2, 3, 4] (expected [2, 3, 5])
stored plate: <NurbsSurface degree=(1,1) net=2x2>
stored lid: open shell with 4 free edges
lid: type=shell, faces=1, edges=4, free_edges=4, ..., closed=False, ...
  ISSUES: shell is open (4 free edges) - not a solid
a: type=curve, degree=2, control_points=3, rational=True, length=3.141592654, ...
solved in 5 iterations: r=1.989436805, h=3.999999985 (max relative target error
  4.63e-09); recipe executed
PASS: area = 50.00000023 (expected 50)
workspace: ['a', 'box', 'cyl', 'gen', 'lid', 'part', 'plate', 'rod']
\end{lstlisting}

The \code{Workspace} is a plain data class. Its \code{objects} attribute is a dictionary, and
there are two methods:
\begin{itemize}
  \item \code{put(name, obj)} stores an object;
  \item \code{get(name)} returns it, or raises \code{KeyError} listing the known names.
\end{itemize}
Objects created in Python can be put into the workspace before a request, and the model's
results can be taken out afterwards. All of them are ordinary \code{NurbsCurve},
\code{NurbsSurface}, \code{TrimmedSurface} and \code{Shell} objects.

\section{Verification}\label{sec:llm-verify}

Language models are good at choosing operations, and poor at noticing that a result is wrong:
a shell that did not close, a radius off by a factor of two, a part in the wrong place. \NF{}
gives the model \emph{measurements} instead of letting it believe its own plan.

\subsection{The tools \texttt{inspect} and \texttt{check}}

\toolref{inspect} reports the measured properties of an object:
\begin{itemize}
  \item \textbf{all objects:} the type and a bounding box from dense samples;
  \item \textbf{curves:} degree, number of control points, whether the curve is rational,
    length, start and end points, and whether it is closed;
  \item \textbf{surfaces and trimmed surfaces:} degree, control net size, area, and the number
    of holes;
  \item \textbf{shells:} faces, edges, free and non-manifold edges, whether the shell is closed
    and consistently oriented, area, and the volume if the shell is closed.
\end{itemize}
Problems are flagged explicitly after \code{ISSUES:}: a degenerate curve or surface (zero
length or area), an empty shell, an open shell with its number of free edges, non-manifold
edges, inconsistent orientation, or a negative volume (a shell oriented inwards).

\toolref{check} compares one measured property with an expected value and answers \code{PASS}
or \code{FAIL} with the measured value. Properties are addressed by path: \code{volume},
\code{area}, \code{length}, \code{closed}, \code{faces}, \code{bbox.min}, \code{bbox.max},
\code{bbox.max[2]}, \code{start}, \dots{}
\begin{itemize}
  \item Numbers and lists of numbers pass within the given absolute \code{tolerance}. By
    default the tolerance is $10^{-6}$ relative, and at least $10^{-9}$ absolute.
  \item Booleans and strings must be equal.
\end{itemize}

\begin{note}[Bounding boxes are sampled]
Bounding boxes are computed from samples, a grid of points on every face. For trimmed faces
only the samples inside the trim region are used. They are therefore slightly inside the true
box for curved or trimmed geometry. In the example, the box of the open \code{lid} runs from
$0.083$ to $0.917$ instead of from 0 to 1. Use a tolerance when checking bounding boxes of such
objects, or check \code{area} and \code{volume} instead, which are integrated accurately.
\end{note}

\subsection{Verification rounds}

After the model's final answer, \code{GeometryAgent.ask} starts up to \code{verify_rounds}
verification rounds:
\begin{enumerate}
  \item The objects created or changed by this request are measured, and a report like that of
    \code{inspect} is sent to the model, together with the instruction to compare it with the
    request.
  \item The model either repairs the geometry with further tool calls, or confirms with an
    answer beginning with \code{VERIFIED}.
  \item If the model made tool calls, the next round measures the objects again. The rounds
    stop as soon as the model answers without calling a tool.
\end{enumerate}
The default is one round. \code{verify_rounds=0} disables verification. The last report is
kept in \code{agent.verification}.

The following script shows this loop without Ollama, using a scripted stand-in for the model.
In its first answer the ``model'' only creates a flat plate. The verification report shows
this, and the model adds a closed cube and checks its volume. A second round confirms the
result.

\example{16_llm_agent_scripted.py}

\begin{lstlisting}[style=text]
answer: VERIFIED: closed unit cube 'cube', volume 1.
  [make_rectangle] -> stored floor: <NurbsSurface degree=(1,1) net=2x2>
  [solid] -> stored cube: closed shell, volume 1
  [check] -> PASS: volume = 1 (expected 1)
last verification report sent to the model:
floor: type=surface, degree=[1, 1], control_points=[2, 2], area=1,
  bbox={min: [0, 0, 0], max: [1, 1, 0]}
cube: type=shell, faces=6, edges=12, free_edges=0, non_manifold_edges=0, closed=True,
  consistently_oriented=True, area=6, volume=1, bbox={min: [0, 0, 0], max: [1, 1, 1]}
workspace: ['cube', 'floor']
\end{lstlisting}

\subsection{The verification functions in Python}

The measuring code is in \code{nurbforge.llm.verify} and works without a model:
\begin{center}
\begin{tabularx}{\textwidth}{@{}lX@{}}
\toprule
Function & Result \\
\midrule
\code{inspect_object(obj)} & Dictionary of measured properties, as listed above. \\
\code{issues(props)} & List of problems found in such a dictionary. \\
\code{describe(name, obj)} & One line: properties and issues. A failing measurement is itself
  reported. \\
\code{report(objects)} & Report for a dictionary of named objects, one line per object. \\
\code{check(obj, prop, expected, tolerance=None)} & \code{"PASS: ..."} or \code{"FAIL: ..."},
  as with the tool. \\
\bottomrule
\end{tabularx}
\end{center}
These functions are handy in your own scripts too. One use is a quick plausibility report
after a long chain of Booleans.

\section{Constraint solving}\label{sec:llm-solve}

Requests often fix a dimension only indirectly: ``a cylinder of height 4 with a lateral area of
50'', ``a lens of volume 2\,cm$^3$'', ``a box whose top is at $z=10$''. Models guess such
values badly. The tool \toolref{solve} lets them state the constraints instead, and computes
the dimensions.

\subsection{The three parts of a \texttt{solve} call}

\begin{lstlisting}[style=json]
{"recipe": [
   {"tool": "make_line", "args": {"name": "gen", "p0": ["$r", 0, 0], "p1": ["$r", 0, "$h"]}},
   {"tool": "revolve", "args": {"name": "cyl", "profile": "gen", "axis_point": [0, 0, 0],
                                "axis_dir": [0, 0, 1], "angle_deg": 360}}],
 "unknowns": {"r": {"start": 1, "min": 0.01, "max": 10},
              "h": {"start": 1, "min": 0.01, "max": 10}},
 "targets": [{"name": "cyl", "property": "area", "value": 50},
             {"name": "cyl", "property": "bbox.max[2]", "value": 4}]}
\end{lstlisting}

\begin{description}[leftmargin=1.5em, style=nextline]
  \item[recipe] A list of ordinary tool calls \code{{"tool": ..., "args": ...}}. Any numeric
    argument, also inside lists, may be written as:
    \begin{itemize}
      \item \code{"$name"}, an unknown;
      \item \code{"${expression}"}, an arithmetic expression over the unknowns.
    \end{itemize}
    Expressions may contain numbers, unknowns, \code{+ - * / **}, parentheses, \code{pi},
    \code{sqrt}, \code{abs}, \code{min}, \code{max}, and \code{sin}, \code{cos}, \code{tan},
    \code{asin}, \code{acos}, \code{atan} in degrees. Example: \code{"${2*r + 1}"}.
  \item[unknowns] Maps each unknown to \code{{"start": ..., "min": ..., "max": ...}}. The bounds
    are optional but strongly recommended. They keep the iteration away from degenerate
    geometry, such as a negative radius.
  \item[targets] Measured properties, as in \code{check}: \code{{"name", "property", "value"}},
    optionally with a \code{weight}. Any property of \code{inspect} can be used.
\end{description}

\subsection{How it works}

\begin{enumerate}
  \item The recipe is executed in a \emph{scratch copy} of the workspace, with the current
    values of the unknowns. The targets are then measured.
  \item The relative target residuals are driven to zero by a bounded Levenberg--Marquardt
    iteration. The Jacobian is computed by forward differences, stepping inwards at the bounds.
    Trial steps that produce invalid geometry are rejected, and the damping is increased.
  \item The solve has converged when the residual norm is below $10^{-8}$; up to $10^{-6}$ is
    accepted. At most 40 iterations are made.
  \item Only after convergence is the recipe executed in the real workspace. Otherwise nothing
    is created, and the model is told the values reached and the remaining error. It can then
    fix the bounds, the recipe or the targets.
\end{enumerate}

{\raggedright For the example above the result is
\code{solved in 5 iterations: r=1.989436805, h=3.999999985}. The exact solution is
$r = 50/(8\pi) = 1.9894368$.

The solver is also available directly, as \code{nurbforge.llm.constraints.solve(ws, recipe,
unknowns, targets, call_tool)}. Expressions can be evaluated with
\code{constraints.evaluate(expr, variables)}.\par}

\section{The command line: \texttt{nurbforge-llm}}

\begin{lstlisting}[style=shell]
nurbforge-llm --model qwen2.5:14b --out model.msh --mesh-size 0.5
\end{lstlisting}

The command starts an interactive loop:
\begin{itemize}
  \item every line typed is passed to \code{agent.ask};
  \item after each answer, the last tool calls are printed with their results;
  \item an empty line (or end of input) quits.
\end{itemize}
With \code{--out}, all surfaces in the workspace are meshed with gmsh on exit and written to
that file. \Cref{app:cli} lists all options.

\section{A session with a real model}

\example{16_llm_ollama.py}

The script needs a running \code{ollama serve}, so it is not executed by the test suite. A
typical model first creates the box and the cylinder with \toolref{solid}, then subtracts them
with \toolref{boolean}, and checks the volume with \toolref{check} or in the verification
round. The expected value is $1000-4\pi\cdot5\approx937.17$.

\section{Adding a tool}

Tools are defined in \code{python/nurbforge/llm/tools.py}, inside \code{_build_tools()}. Each
one is a function decorated with \code{@add(name, description, properties, required)}:

\begin{lstlisting}[style=python]
@add("make_line", "Create a straight line segment.",
     {"name": _NAME, "p0": _VEC, "p1": _VEC}, ["name", "p0", "p1"])
def _(ws, name, p0, p1):
    return ws.put(name, nf.make_line(p0, p1))
\end{lstlisting}

\begin{itemize}
  \item \code{properties} is a JSON schema per argument. \code{_VEC} is a list of three numbers,
    and \code{_NAME} is an object name.
  \item The function receives the workspace and the arguments, and returns a short text.
  \item Exceptions are caught by \code{call_tool} and returned as \code{ERROR:} messages, so
    the function may simply raise.
  \item Keep descriptions short and precise. Small models depend on them.
  \item Convert angles from degrees.
  \item Add a test in \code{tests/python/test_llm.py}, using the scripted \code{FakeClient}.
\end{itemize}
